\unitopen{21}
\begin{candos}
\item opisać rzecz lub osobę przymiotnikiem przed rzeczownikiem (\de{ein neuer Laptop});
\item dobrać końcówkę po rodzajniku określonym i nieokreślonym w mianowniku i bierniku;
\item złożyć reklamację i opisać problem z zamówieniem;
\item porównać dwie oferty z użyciem przymiotników.
\end{candos}

\section{Słownictwo}
\vocabtable{21}

\section{Końcówki przymiotników}
Przymiotnik w roli orzecznika (\de{Der Laptop ist neu}) nie ma końcówki. Przed rzeczownikiem dostaje ją, a jej wybór zależy
od \emph{rodzajnika}. Dwie tabele wystarczą na A2.

\begin{center}\small
\begin{tabular}{@{}lllll@{}}
\toprule
\multicolumn{5}{@{}l}{\textsf{\textbf{Po rodzajniku określonym (\de{der, die, das}, \de{dieser})}}}\\
 & \textsf{\textbf{m}} & \textsf{\textbf{f}} & \textsf{\textbf{n}} & \textsf{\textbf{mnoga}}\\
\midrule
Nominativ & \de{der neu\textbf{e} Laptop} & \de{die neu\textbf{e} Lampe} & \de{das neu\textbf{e} Auto} & \de{die neu\textbf{en} Laptops}\\
Akkusativ & \de{den neu\textbf{en} Laptop} & \de{die neu\textbf{e} Lampe} & \de{das neu\textbf{e} Auto} & \de{die neu\textbf{en} Laptops}\\
\midrule
\multicolumn{5}{@{}l}{\textsf{\textbf{Po rodzajniku nieokreślonym (\de{ein, kein, mein})}}}\\
Nominativ & \de{ein neu\textbf{er} Laptop} & \de{eine neu\textbf{e} Lampe} & \de{ein neu\textbf{es} Auto} & \de{keine neu\textbf{en} Laptops}\\
Akkusativ & \de{einen neu\textbf{en} Laptop} & \de{eine neu\textbf{e} Lampe} & \de{ein neu\textbf{es} Auto} & \de{keine neu\textbf{en} Laptops}\\
\bottomrule
\end{tabular}
\end{center}
Zasada: gdy rodzajnik \emph{nie pokazuje} rodzaju i przypadku (\de{ein} w mianowniku męskim i nijakim), przymiotnik robi to za niego
(\de{-er}, \de{-es}). Po \de{der/die/das}, które już to pokazują, przymiotnik dostaje \de{-e} (m, f, n w mianowniku;
f i n w bierniku) albo \de{-en} (wszystko inne).

\begin{falle}
Polski przymiotnik zgadza się z rzeczownikiem pełną odmianą, więc intuicja jest właściwa, ale końcówki
niemieckie są mniej liczne i mylą się przy rodzaju męskim w bierniku: \de{Ich kaufe einen neu\textbf{en} Laptop}
(nie \de{neuen} bez \de{einen}). Zwróć uwagę: \de{ein} bez końcówki rodzajnika w mianowniku męskim i nijakim wymaga \de{-er} i \de{-es}
na przymiotniku. Przymiotnik na \de{-el}, \de{-er} traci \de{-e-}: \de{teuer $\rightarrow$ ein teurer Laptop}.
\end{falle}

\begin{szwajcaria}
W sklepach internetowych i na szwajcarskich stronach zobaczysz \de{inkl.\ MwSt.} (z VAT, w CH \de{MWST}),
\de{Lieferzeit}, \de{Gratis-Versand}. Prawo zwrotu przy zakupach online bywa krótsze niż w UE
\verify{(w Szwajcarii nie ma ustawowego prawa odstąpienia dla większości zakupów online, sprawdź regulamin sklepu)}.
\end{szwajcaria}

\begin{praca}
Opis sprzętu: \de{Wir brauchen einen schnellen Server und einen stabilen Router.} \de{Das ist ein praktisches Tool.}
\de{Die neue Version hat einen leichteren Aufbau.} W e-mailu: \de{Die gelieferte Ware ist beschädigt.}
\end{praca}

\section{Dialogi}
\textbf{Dialog 1 --- reklamacja}
\begin{dialog}
\turn{A}{Guten Tag, ich habe vor einer Woche einen neuen Bildschirm bestellt.}{Dzień dobry, tydzień temu zamówiłem nowy monitor.}
\turn{B}{Ja, was ist das Problem?}{Tak, jaki jest problem?}
\turn{A}{Der Bildschirm ist beschädigt. Ich möchte ihn zurückgeben und einen anderen bekommen.}{Monitor jest uszkodzony. Chcę go zwrócić i dostać inny.}
\turn{B}{Das tut mir leid. Wir schicken Ihnen einen neuen Bildschirm.}{Przykro mi. Wyślemy Panu nowy monitor.}
\end{dialog}

\section{Ćwiczenia}
\begin{cwiczenia}
\exitem{Uzupełnij: \de{der neu\dots{} Laptop} (Nominativ).}{\de{der neue Laptop}}
\exitem{Uzupełnij: \de{ein neu\dots{} Laptop} (Nominativ).}{\de{ein neuer Laptop}}
\exitem{Uzupełnij: \de{Ich kaufe einen neu\dots{} Laptop.}}{\de{einen neuen Laptop}}
\exitem{Uzupełnij: \de{eine praktisch\dots{} Tasche} (Nominativ).}{\de{eine praktische Tasche}}
\exitem{Uzupełnij: \de{das groß\dots{} Büro} (Nominativ).}{\de{das große Büro}}
\exitem{Uzupełnij: \de{ein groß\dots{} Büro} (Nominativ).}{\de{ein großes Büro}}
\exitem{Uzupełnij: \de{die neu\dots{} Laptops} (Nominativ plural).}{\de{die neuen Laptops}}
\exitem{Uzupełnij: \de{Wir suchen einen stabil\dots{} Server.}}{\de{einen stabilen Server}}
\exitem{Popraw: \de{Ich kaufe ein neuen Laptop.}}{\de{Ich kaufe einen neuen Laptop.}}
\exitem{Popraw: \de{Das ist ein teuerer Laptop.}}{\de{Das ist ein teurer Laptop.}}
\exitem{Przetłumacz: „To jest praktyczny program.”}{\de{Das ist ein praktisches Programm.}}
\exitem{Przetłumacz: „Szukam nowego mieszkania z dużym balkonem.” (\de{mit} + celownik, wzór z jednostki 13)}{\de{Ich suche eine neue Wohnung mit einem großen Balkon.}}
\end{cwiczenia}

\weekend{Opisz trzy rzeczy ze swojego biura lub domu przymiotnikami (po trzy przymiotniki). Poproś AI,
by zadawała pytania: \de{Was für ein Laptop ist das?} Odpowiadaj pełnymi zdaniami.}
